% SOURCE_AUTHORITY: sga5_fr_workpass.tex, lines 8980--9010 (LF-normalized unit).
% SOURCE_SHA256: 791F4EFFC5E02832D5D77ED03518C8156D6F07E4C8238B03545DB93D883FBB28
\subsubsection*{3.2}
Sean \(\ell\) un número primo \(\ne\car k\) y \(X\) un esquema liso de tipo finito sobre \(k\). Si \(Z\) es un ciclo algebraico de codimensión \(d\) de \(X\), en la Exposición IV\footnote{Véase (SGA \(4^{1/2}\) Cycle).} se le ha asociado, para todo entero \(n\), una clase de cohomología
\[
\gamma_X^{(n)}(Z)\in H^{2d}\bigl(X,(\mu_X^{(n)})^{\otimes d}\bigr),
\]
donde \(\mu_X^{(n)}\) denota el \((\mathbb Z/\ell^{\,n+1}\mathbb Z)_X\)-módulo localmente libre de las raíces \(\ell^{\,n+1}\)-ésimas de la unidad. Por otra parte, se ha visto que, si \(n'\ge n\), \(\gamma_X^{(n)}(Z)\) es la imagen de \(\gamma_X^{(n')}(Z)\) en
\[
H^{2d}\bigl(X,(\mu_X^{(n)})^{\otimes d}\bigr)
\]
por el morfismo
\[
H^{2d}\bigl(X,(\mu_X^{(n')})^{\otimes d}\bigr)
\longrightarrow
H^{2d}\bigl(X,(\mu_X^{(n)})^{\otimes d}\bigr)
\]
deducido de la elevación a la potencia \(\ell^{\,n'-n}\):
\[
\mu_X^{(n')}\longrightarrow \mu_X^{(n)}.
\]
Se asocia así a \(Z\) un elemento, denotado por \(\gamma_X(Z)\):
\[
\gamma_X(Z)\in H^{2d}(X,\mathbb Z_\ell(d)),
\]
donde \(\mathbb Z_\ell(d)\) se define en (1.3.4). Más generalmente, si \(Z\) es un ciclo de \(X\), se le asocia por linealidad una clase de cohomología \(\ell\)-ádica en
\[
\bigoplus_d H^{2d}(X,\mathbb Z_\ell(d)),
\]
que es, de hecho, el límite proyectivo de las clases de cohomología asociadas a \(Z\) en los
\[
\bigoplus_d H^{2d}\bigl(X,(\mu_X^{(n)})^{\otimes d}\bigr).
\]
